\documentclass[10pt]{article}
\usepackage[margin=1in]{geometry}
\usepackage[T1]{fontenc}
\usepackage{lmodern,microtype,amsmath,amssymb,amsthm}
\theoremstyle{definition}
\newtheorem{assumption}{Assumption}
\begin{document}
\section*{Generalized Smooth Finite-Sum Minimization}
We consider minimizing $F(x)=n^{-1}\sum_{i=1}^n f_i(x)$ on a real Hilbert space $E$, where $n\geq1$ is an integer. Each component $f_i:E\to\mathbb{R}$ is convex, continuously differentiable, and $(L_0,L_1)$-smooth, with $L_0>0$ and $L_1\geq0$. A global minimizer $x^\star$ satisfies bounds on its distance from a given starting point $x^0\in E$, the initial objective gap, and the average component-gradient norm at $x^\star$. At a query $(i,x)$, the oracle returns $(f_i(x),\nabla f_i(x))$. For $\varepsilon>0$, the goal is a measurable output $\bar{x}\in E$ with integrable objective gap satisfying $\mathbb{E}_{P}[F(\bar{x})-F(x^\star)]\leq\varepsilon$, where $P$ is the law of the algorithm's randomness.

The optimization problem is
\[
\min_{x\in E} F(x).
\]
The norm $\|\cdot\|$ is induced by the inner product on $E$, which is equipped with a measurable structure. The algorithm receives $E$, $n$, $L_0$, $L_1$, $G\geq0$, $R>0$, $\Delta>0$, $\varepsilon>0$, and $x^0$.

Our assumptions:
\begin{assumption}[Component convexity]
For every $i\in\{1,\ldots,n\}$, $x,y\in E$, and $t\in[0,1]$,
\[
f_i((1-t)x+ty)\leq (1-t)f_i(x)+t f_i(y).
\]
\end{assumption}

\begin{assumption}[Generalized smoothness]
Each $f_i$ is Fr\'echet differentiable everywhere, and $\nabla f_i:E\to E$ is continuous. For every $i\in\{1,\ldots,n\}$, $x,y\in E$, and $M\in\mathbb{R}$ satisfying
\[
\|\nabla f_i(x+\theta(y-x))\|\leq M
\qquad\text{for all }\theta\in[0,1],
\]
we have
\[
\|\nabla f_i(x)-\nabla f_i(y)\|
\leq (L_0+L_1M)\|x-y\|.
\]
\end{assumption}

\begin{assumption}[Minimizer and initial bounds]
There exists $x^\star\in E$ such that
\[
\begin{aligned}
F(x^\star)&\leq F(x) &&\text{for every }x\in E,\\
\|x^0-x^\star\|&\leq R,\\
F(x^0)-F(x^\star)&\leq\Delta,\\
\frac{1}{n}\sum_{i=1}^n\|\nabla f_i(x^\star)\|&\leq G.
\end{aligned}
\]
\end{assumption}

The algorithm may query any component at any point in $E$; each query returns the exact component value and gradient together. After the allotted oracle calls, its output $\bar{x}$ must be measurable, its objective gap must be integrable under $P$, and it must satisfy
\[
\mathbb{E}_{P}\bigl[F(\bar{x})-F(x^\star)\bigr]\leq\varepsilon.
\]
\end{document}
